\documentclass[11pt, a4paper]{article}

% ====================================
% 引入刚才写的宏包
% ====================================
\usepackage{optstyle} 
\usepackage{amsmath}
\fancyhead[R]{\sffamily\footnotesize Fisher Markets}
% ====================================
% 笔记信息
% ====================================
\title{\textbf{\Huge Fisher Markets}\\\Large Convex Programming and Duality}
\author{Zhiyun Guo}
\date{January 28, 2026}

\begin{document}

\maketitle

% ====================================
% 正文开始
% ====================================

\section{Preliminaries}

\subsection{Notation}
\begin{itemize}
    \item $M$: a set of $m$ divisible goods
    \item $N$: a set of $n$ buyers
    \item $B_i$: Each buyer i's budget (fixed)
    \item $p_j\in \mathbb{R}_+$: price of good $j$ 
    \item $x_{ij}\in\mathbb{R}_+$: amount of good $j$ allocated to buyer $i$
\end{itemize}
\paragraph{Standing assumptions.} Budgets satisfy $B_i>0$. Goods are divisible and have unit supply after rescaling positive supplies. For the linear-market programs below, each buyer values some good, and each good is valued by some buyer. Nonnegativity is allowed for $v_{ij}$; zero-valued pairs are excluded when taking logarithms or ratios.
\subsection{Utility}
\begin{enumerate}
    \item \textit{Linear}: there are parameters $v_{ij}\in\mathbb{R}_+$ specifying the value derived by buyer $i$ from one unit of good $j$ which guarantees \textbf{Additivity}. i.e., $u_i(x)=\sum_{j\in M}v_{ij}x_{ij}$.
    \item\textit{Quasi-linear}: agents have utility for the money spent as well, i.e., $u_i(x)=\sum_{j\in M}(v_{ij}-p_j)x_{ij}$.
    \item \textit{Leontief}: given coefficients $a_{ij}\ge0$, with at least one positive coefficient per buyer, $u_i(x)=\min_{j:a_{ij}>0}x_{ij}/a_{ij}$. Goods with $a_{ij}=0$ are not required inputs.
    \item \textit{Constant elasticity of substitution (CES)}: for example, when $0<\rho<1$,
\[
u_i(x_i)=\left(\sum_j\alpha_{ij}x_{ij}^{\rho}\right)^{1/\rho},
\]
with $\alpha_{ij}\ge0$ and some positive coefficient. Other parameter regimes, including complementary goods and the Cobb--Douglas limit, require the corresponding domain/limit conventions. 
\end{enumerate}
\section{Market equilibrium}
An \textit{equilibrium} is formed if the following conditions hold.
\begin{enumerate}
    \item The allocation of each buyer $i$ \textbf{maximizes his utility}, subject to the budget constraint $\sum_j p_j x_{ij}\leq B_i$
    \item Each good $j$ that has a price $p_j>0$ is \textbf{allocated fully}, i.e., $\sum_i x_{ij}=1$. A good is allowed to have $p_j=0$ as long as $\sum_i x_{ij}\leq 1$.
\end{enumerate}
\subsection{Generalization of Fisher's model}
\subsubsection{Spending-Restricted (SR) model}
\begin{itemize}
    \item Each seller $j$ has an upper bound on the amount of money he wants to earn: $c_j$
    \item Once he earns $c_j$, selling the least amount of his good, he wants to take back the unsold portion of his good, i.e., the amount of money spent on the good of seller $j$ is restricted by $c_j$
    \item In equilibrium, the second equilibrium condition above is modified to $\forall j\in M,\sum_i x_{ij}\leq 1,\sum_i p_j x_{ij}\leq c_j$, and either 
    \[\sum_i x_{ij} = 1 \text{ or }\sum_i p_j x_{ij}=c_j, \text{ or } p_j=0. \]
\end{itemize}
\subsubsection{Utility-Restricted (UR) model}
\begin{itemize}
    \item Each buyer $i$ has an upper bound on the utility they want to derive in the market: $d_i$
    \item Once buyer $i$ derives utility $d_i$, spending the least amount of money at prices $p$, she wants to keep the left-over money. i.e., the utility of buyer $i$ is restricted by $d_i$
    \item The first equilibrium condition is modified to $\forall i\in N,\; u_i(x)\leq d_i, \text{ and }\sum_j p_j x_{ij}\leq B_i$, and either 
    \[x\text{ minimizes }\sum_j p_j x_{ij}\; s.t.\; u_i(x)=d_i,\text{ or maximizes }u_i(x)\; s.t. \sum_j p_j x_{ij}\leq B_i.\]
\end{itemize}
\section{Convex programming duality}
\subsection{Fenchel Conjugate}
\begin{definition}{Fenchel conjugate}{fj}
For $f:\mathbb R^n\to\mathbb R\cup\{+\infty\}$, its conjugate is
\[
f^*(\mu)=\sup_{x\in\operatorname{dom}f}\{\mu^\top x-f(x)\}.
\]
The conjugate may be extended-valued. It is convex as a supremum of affine functions; strict convexity and differentiability do not follow in general.
\end{definition}

\noindent\textbf{Useful rules and assumptions.}
\begin{itemize}
\item If $f$ is proper, closed and convex, $f^{**}=f$ (Fenchel--Moreau).
\item If $f(x)=\sum_i f_i(x_i)$ on a product domain, $f^*(\mu)=\sum_i f_i^*(\mu_i)$.
\item If $g=cf$ with $c>0$, then $g^*(\mu)=cf^*(\mu/c)$.
\item If $g(x)=f(cx)$ with scalar $c\ne0$, then $g^*(\mu)=f^*(\mu/c)$.
\item If $g(x)=f(x+a)$, then $g^*(\mu)=f^*(\mu)-\mu^\top a$.
\item For proper closed convex $f$, Fenchel equality is equivalent to a subgradient pair:
\[
f(x)+f^*(\mu)=\mu^\top x
\quad\Longleftrightarrow\quad
\mu\in\partial f(x)
\quad\Longleftrightarrow\quad
x\in\partial f^*(\mu).
\]
Gradient identities require differentiability at the corresponding points.
\end{itemize}
For example, the conjugate of $f(x)=0$ on $\mathbb R$ is $0$ at $\mu=0$ and $+\infty$ elsewhere, illustrating why general differentiability cannot be asserted.

\noindent\textbf{Examples.}
\begin{itemize}
\item $f(x)=x^2/2$ on $\mathbb R$: $f^*(\mu)=\mu^2/2$.
\item $f(x)=-\log x$ for $x>0$: $f^*(\mu)=-1-\log(-\mu)$ for $\mu<0$, and $+\infty$ otherwise.
\item $f(x)=x\log x$ on $x\ge0$, with $0\log0=0$: the optimizer for finite $\mu$ is $x=e^{\mu-1}$, hence $f^*(\mu)=e^{\mu-1}$.
\end{itemize}
See Boyd and Vandenberghe, Section 3.3, for conjugacy rules and domains \cite{boyd2004}.

\subsection{Convex Programs with linear constraints}

\noindent
\begin{minipage}[t]{0.48\textwidth}
    \textbf{Primal (Maximization):}
    \begin{align*}
        \max \quad & \sum_{i} c_i x_i - f(x) \\
        \text{s.t.} \quad & \sum_{i} a_{ij} x_i \leq b_j, \quad \forall j \\
                          & x_i \ge 0, \quad \forall i
    \end{align*}
\end{minipage}
\hfill \vrule{} \hfill
\begin{minipage}[t]{0.48\textwidth}
    \textbf{Dual (Minimization):}
    \begin{align*}
        \min \quad & \sum_{j} b_j \lambda_j + f^*(\mu) \\
        \text{s.t.} \quad & \sum_{j} a_{ij} \lambda_j \ge c_i - \mu_i, \quad \forall i \\
                          & \lambda_j \geq 0, \quad \forall j
    \end{align*}
\end{minipage}

\vspace{1em}
\begin{interpretation}
    \textbf{Strong Duality \& Slackness.}
    Under a suitable Slater constraint qualification and finite attained optima, primal and dual values agree. The solution is characterized by:
    \begin{enumerate}
        \item \textbf{Complementary Slackness:} 
        \[ x_i > 0 \implies \text{Dual Constraint is Tight}, \quad \lambda_j > 0 \implies \text{Primal Constraint is Tight}. \]
        \item \textbf{Conjugate Pair:} $\mu\in\partial f(x)$ and $x\in\partial f^*(\mu)$; write $\mu=\nabla f(x)$ only when $f$ is differentiable at $x$.
    \end{enumerate}
\end{interpretation}
\noindent Alternatively, for a minimization primal, the forms are reversed:

\noindent
\begin{minipage}[t]{0.48\textwidth}
    \textbf{Primal (Minimization):}
    \begin{align*}
        \min \quad & \sum_{i} c_i x_i + f(x) \\
        \text{s.t.} \quad & \sum_{i} a_{ij}x_i \ge b_j, \quad \forall j \\
                          & x_i \ge 0.
    \end{align*} 
\end{minipage}
\hfill \vrule{} \hfill
\begin{minipage}[t]{0.48\textwidth}
    \textbf{Dual (Maximization):}
    \begin{align*}
        \max \quad & \sum_{j} b_j \lambda_j - f^*(\mu) \\
        \text{s.t.} \quad & \sum_{j} a_{ij}\lambda_j \le c_i + \mu_i, \quad \forall i \\
                          & \lambda_j \ge 0.
    \end{align*}
\end{minipage}

\section{Fisher Markets: Eisenberg-Gale \& Shmyrev}

The Eisenberg-Gale (EG) program is fundamental in algorithmic game theory as it captures the equilibrium of linear Fisher markets, which maximizes budget-weighted log utilities, equivalently the weighted Nash product $\prod_i u_i^{B_i}$. Its dual variables provide the equilibrium prices.

\subsection{The Eisenberg-Gale (EG) Program}

\begin{lemma}{EG Duality}{eg_dual}
    The EG Primal and its Dual are related as follows. The dual variables $p_j$ correspond to item prices.
\end{lemma}

\noindent
\begin{minipage}[t]{0.48\textwidth}
    \textbf{Primal (EG Program):}
    \begin{align*}
        \max \quad & \sum_{i} B_i \log u_i \\
        \text{s.t.} \quad & u_i \leq \sum_{j} v_{ij} x_{ij}, \quad \forall i \\
        & \sum_{i} x_{ij} \leq 1, \quad \forall j \\
        & x_{ij} \geq 0.
    \end{align*}
\end{minipage}
\hfill \vrule{} \hfill
\begin{minipage}[t]{0.48\textwidth}
    \textbf{Dual (Eq. 1):}
    \begin{align*}
        \min \quad & \sum_{j} p_j - \sum_{i} B_i \log(\beta_i) \\
        \text{s.t.} \quad & p_j \geq v_{ij} \beta_i, \quad \forall i, j \\
        & p_j\ge0,\quad\beta_i>0.
    \end{align*}
\end{minipage}

\begin{remark}
    The displayed dual drops the additive constant $C_B=\sum_i(B_i\log B_i-B_i)$. At an optimum, $\beta_i=\min_{j:v_{ij}>0}p_j/v_{ij}$; this is the reciprocal of buyer $i$'s maximum bang-per-buck.
\end{remark}

\begin{proof}[Derivation and equilibrium interpretation]
Introduce multipliers $\beta_i\ge0$ for $u_i\le\sum_jv_{ij}x_{ij}$ and $p_j\ge0$ for supply. Maximizing the Lagrangian over $x_{ij}\ge0$ is finite precisely when $p_j\ge\beta_i v_{ij}$. Maximizing over $u_i>0$ gives
\[
\sup_{u_i>0}(B_i\log u_i-\beta_i u_i)
=B_i\log B_i-B_i-B_i\log\beta_i,
\]
with $\beta_i>0$. This yields the displayed dual plus $C_B$. At an optimum, $u_i=B_i/\beta_i$ and complementary slackness gives $x_{ij}>0\Rightarrow p_j=\beta_i v_{ij}$. Thus each buyer uses only maximum bang-per-buck goods and spends $\sum_jp_jx_{ij}=\beta_i u_i=B_i$. Supply complementarity supplies market clearing.
\end{proof}

\subsection{Connection to Shmyrev's Program}

We can derive the Shmyrev Convex Program by applying a change of variables to the EG Dual.

\begin{interpretation}
    \textbf{Logarithmic transformation and domains.} Under the standing assumptions, equilibrium prices and $\beta_i$ are positive. Only pairs with $v_{ij}>0$ enter the logarithmic constraints; in the bidding program set $b_{ij}=0$ on zero-valued pairs.
    To linearize the constraints in the EG Dual, we define:
    \[ q_j = \log p_j \quad \text{and} \quad \gamma_i = -\log \beta_i \]
    The constraint $p_j \geq v_{ij} \beta_i$ transforms into $\log p_j \geq \log v_{ij} + \log \beta_i$, which becomes:
    \[ \gamma_i + q_j \geq \log v_{ij} \]
\end{interpretation}
\noindent This transformation leads to a new pair of dual programs:

\noindent
\begin{minipage}[t]{0.48\textwidth}
    \textbf{Transformed Dual (Eq. 2):}
    \begin{align*}
        \min \quad & \sum_{j} e^{q_j} + \sum_{i} B_i \gamma_i \\
        \text{s.t.} \quad & \gamma_i + q_j \geq \log v_{ij}, \quad \forall i, j.
    \end{align*}
\end{minipage}
\hfill \vrule{} \hfill
\begin{minipage}[t]{0.48\textwidth}
    \textbf{Shmyrev Program (CP):}
    \begin{align*}
        \max \quad & \sum_{i,j} b_{ij} \log v_{ij} - \sum_{j} (p_j \log p_j - p_j) \\
        \text{s.t.} \quad & \sum_{i} b_{ij} = p_j, \quad \forall j \\
        & \sum_{j} b_{ij} = B_i, \quad \forall i \\
        & b_{ij} \geq 0.
    \end{align*}
\end{minipage}

\vspace{1em}
\noindent The Shmyrev objective includes $+\sum_jp_j$; this is constant because $\sum_jp_j=\sum_iB_i$ on its feasible set. Consequently dropping this term leaves the same maximizers. The bids recover allocations by $x_{ij}=b_{ij}/p_j$ for positive prices.

\subsection{Quasi-Linear Utilities}

For markets where buyers have \textbf{quasi-linear utilities} (utility form: $\sum_j (v_{ij} - p_j)x_{ij}$), the standard EG Primal is hard to generalize directly. However, the \textit{Dual} generalizes naturally.

\begin{lemma}{Quasi-Linear Duality}{ql_dual}
    By adding the constraint $\beta_i \leq 1$ to the EG Dual, we capture quasi-linear utilities. The corresponding primal is then derived.
\end{lemma}

\noindent
\begin{minipage}[t]{0.48\textwidth}
    \textbf{Dual (Modified EG Type):}
    \begin{align*}
        \min \quad & \sum_{j} p_j - \sum_{i} B_i \log(\beta_i) \\
        \text{s.t.} \quad & p_j \geq v_{ij} \beta_i, \quad \forall i, j \\
        & p_j\ge0,\quad 0<\beta_i\leq1, \quad \forall i.
    \end{align*}
\end{minipage}
\hfill \vrule{} \hfill
\begin{minipage}[t]{0.48\textwidth}
    \textbf{Primal (Derived for Quasi-Linear):}
    \begin{align*}
        \max \quad & \sum_{i} B_i \log u_i - \sum_i m_i \\
        \text{s.t.} \quad & u_i \leq \sum_{j} v_{ij} x_{ij} + m_i, \quad \forall i \\
        & \sum_{i} x_{ij} \leq 1, \quad \forall j \\
        & x_{ij}, m_i \geq 0.
    \end{align*}
\end{minipage}
\noindent Here $m_i$ represents retained money in the utility expression (up to the additive constant $B_i$). The penalty is summed over buyers. As above, the displayed dual omits $C_B$; it is equivalent up to an additive constant, rather than equal in objective value as written.

\begin{thebibliography}{9}
\bibitem{devanur2016} N. R. Devanur, K. Jain, T. Mai, V. V. Vazirani, and S. Yazdanbod. \textit{New Convex Programs for Fisher's Market Model and its Generalizations}. 2016. \href{https://arxiv.org/abs/1603.01257}{Author manuscript}.
\bibitem{boyd2004} S. Boyd and L. Vandenberghe. \textit{Convex Optimization}. Cambridge University Press, 2004. \href{https://web.stanford.edu/~boyd/cvxbook/}{Author's book page}.
\end{thebibliography}
\end{document}
